% Test bed of the template. After changing the template, the packages, or the commands, check
% here that fonts, page layout, acronyms, and figures still work. Remove this section and its
% \input in secs/appendix.tex from your paper.
\section{Test Bed}\label{sec:TestBed}
This appendix tests the fonts, the page layout, the acronyms, and the figures of the template.

\subsection{Fonts}
Each line uses a different font, with the old font commands that some classes still use.

\noindent Waltz, bad nymph, for quick jigs vex.

\noindent{\rm{Waltz, bad nymph, for quick jigs vex.}}

\noindent{\sf{Waltz, bad nymph, for quick jigs vex.}}

\noindent{\tt{Waltz, bad nymph, for quick jigs vex.}}

\noindent{\bf{Waltz, bad nymph, for quick jigs vex.}}

\noindent{\it{Waltz, bad nymph, for quick jigs vex.}}

\noindent{\sl{Waltz, bad nymph, for quick jigs vex.}}

\noindent{\sc{Waltz, bad nymph, for quick jigs vex.}}

\noindent Math: $a$, $\alpha$, $\vect{a}$, $\nvec{a}$, $\pnt{a}$, $\mat{A}$, $\quat{a}$, $a_k$, $a_\mathrm{max}$, $\nreal^3$.

\noindent Numbers: \num{1152}, \num{0.444}, \textbf{\num{0.444}}, \qty{4.5}{\milli\second}, \qty{12}{\percent}.

\subsection{Page Layout}
The text is \the\textwidth{} wide, and a column is \the\columnwidth{} wide.
Draw figures at these widths.
\Cref{fig:TestLayout} shows the layout of a page with all its lengths.

\begin{figure*}[p]
    \centering
    % \layout draws the page at half its size and lists its lengths. Here, it draws one page
    % only, and it neither ends the page nor changes its style.
    {\csname @twosidefalse\endcsname
     \renewcommand{\newpage}{}%
     \renewcommand{\thispagestyle}[1]{}%
     \layout}
    % The table of lengths below the drawing has zero height; make room for it.
    \vspace{80pt}
    \caption{\captiontitle{Page Layout}
        The page at half its size with the lengths of the template.
    }
    \label{fig:TestLayout}
\end{figure*}

\subsection{Acronyms}
% \glsreset makes the next \ac the first use again.
\glsreset{GPU}
The first time, we use a \ac{GPU}.
Next time, we use a \ac{GPU}.
Now, we even have multiple \acp{GPU}.
In headings and captions, write \acs{GPU}.

\subsection{Figures}
\Cref{fig:TestLatexText,fig:TestInkscapeText} include an SVG figure with two options of \texttt{\textbackslash includesvg}:
\begin{itemize}
    \item \texttt{inkscapelatex=true}, the default: \LaTeX{} typesets the text.
    \item \texttt{inkscapelatex=false}: Inkscape renders the text with the fonts of the SVG file.
\end{itemize}
\Cref{fig:TestRaster} includes a PNG file directly.

\begin{figure}[t]
    \centering
    \includesvg[width=\columnwidth, inkscapelatex=true]{figs/PDFTexFigure}
    \caption{\captiontitle{Text Typeset by \LaTeX}
        \LaTeX{} typesets the text of the SVG, so fonts and math match the paper.
    }
    \label{fig:TestLatexText}
\end{figure}

\begin{figure}[t]
    \centering
    \includesvg[width=\columnwidth, inkscapelatex=false]{figs/SVGFigure}
    \caption{\captiontitle{Text Rendered by Inkscape}
        The fonts come from the SVG file.
        Avoid this in papers.
    }
    \label{fig:TestInkscapeText}
\end{figure}

\begin{figure}[t]
    \centering
    \includegraphics[width=0.5\columnwidth]{imgs/checkerboard.png}
    \caption{\captiontitle{Raster Image}
        The texture of \cref{fig:Teaser} with $256\times256$ texels as a PNG file.
    }
    \label{fig:TestRaster}
\end{figure}
